\documentclass[10pt]{article}
\usepackage[letterpaper,margin=0.42in]{geometry}
\usepackage{fontawesome5}
\usepackage{xcolor}
\usepackage{enumitem}
\usepackage{titlesec}
\usepackage[hidelinks]{hyperref}
\usepackage{fancyhdr}

\definecolor{accent}{RGB}{0,100,200}

\titleformat{\section}{\large\bfseries\color{accent}}{}{0em}{}[\titlerule]
\pagestyle{fancy}
\fancyhf{}
\renewcommand{\headrulewidth}{0pt}
\setlist[itemize]{itemsep=1.5pt, topsep=2pt}

\newcommand{\resumeItem}[1]{\item\small{#1}}
\newcommand{\resumeSubheading}[4]{
  \vspace{4pt}
  \noindent
  \begin{tabular*}{\textwidth}[t]{@{}l@{\extracolsep{\fill}}r@{}}
    \textbf{#1} & #2 \\
    \textit{\small #3} & \textit{\small #4} \\
  \end{tabular*}\vspace{-4pt}
}

\begin{document}

\begin{center}
    {\Huge \textbf{Piter Z. Garcia Bautista}} \\
    \vspace{4pt}
  {\large Machine Learning, Quantum Computing, and Reproducible Quantum Systems Research} \\
    \vspace{4pt}
    \faMapMarker* \, Rochester, NY \quad
    \faPhone \, +1 (646) 288-3300 \quad
    \faEnvelope \, \href{mailto:pzg8794@rit.edu}{pzg8794@rit.edu} \\
    \faLinkedin \, \href{https://linkedin.com/in/garciapiterz}{linkedin.com/in/garciapiterz} \quad
    \faGithub \, \href{https://github.com/pzg8794}{github.com/pzg8794} \quad
    \faGlobe \, \href{https://garciapiterz.com}{garciapiterz.com}
\end{center}

\section{Summary}
Graduate researcher with an M.S. in Computer Science and current graduate work in Data Science at Rochester Institute of Technology. Brings industry experience building reliable data and machine learning systems together with current AI and quantum research focused on reproducible evaluation, quantum routing, resource allocation, and learning under stochastic and adversarial conditions. Research interests center on structured machine learning for reliable quantum systems, robust experimentation, and decision-making from noisy and limited observations.

\section{Technical Skills}
\textbf{Programming:} Python, R, Java, C++, C\#, SQL, MATLAB, JavaScript, Bash \\
\textbf{Machine Learning / AI:} PyTorch, TensorFlow, Keras, scikit-learn, XGBoost, contextual bandits, neural bandits, supervised learning, model evaluation, error analysis \\
\textbf{Data Science / Statistics:} pandas, NumPy, SciPy, Matplotlib, Plotly, Seaborn, reproducible benchmarking, experiment design, applied statistics \\
\textbf{Quantum / Research Focus:} quantum routing, qubit allocation, stochastic and adversarial evaluation, multi-run validation, reliability-oriented benchmarking \\
\textbf{Research Tooling:} Git/GitHub, Linux, Docker, LaTeX, AWS, GCP

\section{Relevant Experience}

\resumeSubheading
{Graduate Assistant -- Quantum \& AI Research}{Fall 2025 -- Present}
{Software Engineering Department, Rochester Institute of Technology}{Rochester, NY}
\begin{itemize}[leftmargin=1.2em]
\resumeItem{Build reproducible multi-testbed evaluation workflows for quantum routing and resource-allocation studies across local, Colab, and GCP environments, with shared datasets, structured logging, automated validation, and comparison-ready outputs.}
\resumeItem{Develop Python-based experimental systems for contextual bandits, adversarial neural bandits, and dynamic qubit-allocation analysis under stochastic and non-stationary quantum-network conditions.}
\resumeItem{Produce technical reports and manuscript-ready analyses that compare algorithms across attack scenarios, allocation strategies, and capacity regimes, emphasizing robustness, interpretability, and disciplined evaluation.}
\end{itemize}

\resumeSubheading
{AI Data Solutions Engineer}{Jan 2021 -- Sep 2022}
{VIOME Inc.}{New York, NY}
\begin{itemize}[leftmargin=1.2em]
\resumeItem{Developed applied AI and machine learning workflows in Python for healthcare-oriented use cases, including preprocessing, feature engineering, model experimentation, and evaluation.}
\resumeItem{Worked in AWS-backed environments to support reproducible data preparation and reliable evaluation for predictive and recommendation-oriented applications.}
\end{itemize}

\resumeSubheading
{Data Solutions Engineer}{Sep 2022 -- Aug 2024}
{VEDADATA Inc.}{New York, NY}
\begin{itemize}[leftmargin=1.2em]
\resumeItem{Designed large-scale Python and pandas analytics workflows for ingestion, cleaning, preprocessing, validation, and downstream analysis.}
\resumeItem{Applied data-quality checks and statistical reasoning to improve traceability, reliability, and usability in production data systems.}
\end{itemize}

\section{Education}

\resumeSubheading
{Master of Science in Data Science}{Expected Aug 2026}
{Rochester Institute of Technology}{Rochester, NY}
\begin{itemize}[leftmargin=1.2em]
\resumeItem{Relevant coursework: Data Science Foundations, Software Engineering for Data Science, Data-Driven Knowledge Discovery, Neural Networks, Applied Statistics, Bioinformatics Algorithms.}
\end{itemize}

\resumeSubheading
{Master of Science in Computer Science}{2015}
{Rochester Institute of Technology}{Rochester, NY}
\begin{itemize}[leftmargin=1.2em]
\resumeItem{Concentrations: Artificial Intelligence, Big Data Analytics, Computer Graphics.}
\end{itemize}

\resumeSubheading
{Additional Graduate Study: M.S. in Teaching Computer Science K--12}{Expected Aug 2026}
{University of Rochester, Warner School of Education}{Rochester, NY}
\begin{itemize}[leftmargin=1.2em]
\resumeItem{NSF Robert Noyce Scholar; Advanced Certificate in Leadership in Disability and Inclusive Practices.}
\end{itemize}

\section{Selected Fit for This Role}
\begin{itemize}[leftmargin=*]
\item \small{Current graduate assistant research on quantum routing, dynamic qubit allocation, contextual bandits, and adversarial neural bandits with reproducible multi-run evaluation across noisy and changing conditions.}
\item \small{Strong Python-based machine learning and experiment-design background aligned with learning from limited observations, robust benchmarking, and careful analysis of reliability tradeoffs in quantum settings.}
\item \small{Clear motivation to extend current quantum-systems ML work toward structured learning for reliable quantum computing, including noise-aware modeling and error-mitigation research in the KTH/WASP environment.}
\end{itemize}

\section{References}
\begin{itemize}[leftmargin=*]
\item \small{Professor Travis Desell, Rochester Institute of Technology -- advisor and reference for data science, machine learning, and computational research readiness.}
\item \small{Dr. Daniel Krutz, Rochester Institute of Technology -- reference for technical rigor, software engineering, and research preparation.}
\item \small{Additional reference details available on request.}
\end{itemize}

\end{document}